%%%%%%%%%%%%%%%%%%%%%%%%%%%%%%%%%%%%%%%%%
% University Assignment Title Page 
% LaTeX Template
% Version 1.0 (27/12/12)
%
% This template has been downloaded from:
% http://www.LaTeXTemplates.com
%
% Original author:
% WikiBooks (http://en.wikibooks.org/wiki/LaTeX/Title_Creation)
%
% License:
% CC BY-NC-SA 3.0 (http://creativecommons.org/licenses/by-nc-sa/3.0/)
% 
% Instructions for using this template:
% This title page is capable of being compiled as is. This is not useful for 
% including it in another document. To do this, you have two options: 
%
% 1) Copy/paste everything between \begin{document} and \end{document} 
% starting at \begin{titlepage} and paste this into another LaTeX file where you 
% want your title page.
% OR
% 2) Remove everything outside the \begin{titlepage} and \end{titlepage} and 
% move this file to the same directory as the LaTeX file you wish to add it to. 
% Then add \input{./title_page_1.tex} to your LaTeX file where you want your
% title page.
%
%%%%%%%%%%%%%%%%%%%%%%%%%%%%%%%%%%%%%%%%%
%\title{Title page with logo}
%----------------------------------------------------------------------------------------
%	PACKAGES AND OTHER DOCUMENT CONFIGURATIONS
%----------------------------------------------------------------------------------------

\documentclass[12pt]{article}
\usepackage[english]{babel}
\usepackage[utf8x]{inputenc}
\usepackage{amsmath}
\usepackage{graphicx}
\usepackage[colorinlistoftodos]{todonotes}
\usepackage{pifont}
\usepackage{skull}
\usepackage{wasysym}
\usepackage{skak}

\begin{document}

\begin{titlepage}

\newcommand{\HRule}{\rule{\linewidth}{0.5mm}} % Defines a new command for the horizontal lines, change thickness here

\center % Center everything on the page
 
%----------------------------------------------------------------------------------------
%	HEADING SECTIONS
%----------------------------------------------------------------------------------------

\textsc{\LARGE Università degli studi di Milano-Bicocca}\\[1cm] % Name of your university/college
\textsc{\Large Decision Models}\\[0.3cm] % Major heading such as course name
\textsc{\large Final Project}\\[0.1cm] % Minor heading such as course title

%----------------------------------------------------------------------------------------
%	TITLE SECTION
%----------------------------------------------------------------------------------------

\HRule \\[0.4cm]
{ \huge \bfseries Methodological Approaches for Multi-agent RL games}\\[0.4cm] % Title of your document
\HRule \\[1.5cm]
 
%----------------------------------------------------------------------------------------
%	AUTHOR SECTION
%----------------------------------------------------------------------------------------

\large
\emph{Authors:}\\
Riccardo Cervero - 794126\\   
Federico Moiraghi - 799735\\ 
Pranav Kasela - 846965\\[1cm] 
%----------------------------------------------------------------------------------------
%	DATE SECTION
%----------------------------------------------------------------------------------------

{\large \today}\\[2cm] 

%----------------------------------------------------------------------------------------
%	LOGO SECTION
%----------------------------------------------------------------------------------------

\includegraphics[width=0.2\textwidth]{logo.png}\\[1cm] 
%----------------------------------------------------------------------------------------

\vfill % Fill the rest of the page with whitespace

\end{titlepage}


\begin{abstract}

\end{abstract}

%The introduction should provide a clear statement of the problem posed by the project, and why the problem is of interest. It should reflect the scenario, if available. If needed, the introduction also needs to present background information so that the reader can understand the significance of the problem. A brief summary of the hypotheses and the approach your group used to solve the problem should be given, possibly also including a concise introduction to theory or concepts used later to analyze and to discuss the results.

%Methodological Approach is the central and most important section of the report. Its objective must be to show, with linearity and clarity, the steps that have led to the definition of a decision model. The description of the working hypotheses, confirmed or denied, can be found in this section together with the description of the subsequent refining processes of the models. Comparisons between different models (e.g. heuristics vs. optimal models) in terms of quality of solutions, their explainability and execution times are welcome.
%Do not attempt to describe all the code in the system, and do not include large pieces of code in this section, use pseudo-code where necessary. Complete source code should be provided separately (in Appendixes, as separated material or as a link to an on-line repo). Instead pick out and describe just the pieces of code which, for example: are especially critical to the operation of the system; you feel might be of particular interest to the reader for some reason; illustrate a non-standard or innovative way of implementing an algorithm, data structure, etc..You should also mention any unforeseen problems you encountered when implementing the system and how and to what extent you overcame them. Common problems are: difficulties involving existing software.

\section{Introduction}
Reinforcement Learning is the branch of Machine Learning gathering several techniques designed to create a system, an agent, capable of learning the best way to achieve goals within an environment, whose elements - such as, for example, obstacles or objectives - can also vary over time. In other words, this agent simulate the process of human understanding, by exploring the initially unknown environment and consequently receiving rewards or penalization based on the goodness of the choices made.\\*
In this study, we especially set ourselves the goals of:
\begin{enumerate}
  \item Using different methodologies to perform the same Reinforcement Learning experiment (Game 1);
  \item Deepening the dynamics among multiple separate agent, each of which will learn automatically and independently how to win the game by looking and reacting to other agents' choices. In particular, the two dynamics of "competition" and "cooperation" between agents will be examined.
\end{enumerate}
To do so, four "worlds" have been built, with different rules and elements - such as fear or randomly moving obstacles -, within the agents will be trained over several epochs to maximize the reward, beat the rivals or help the allies in the best way possible.\\*
The following sections will present games' and agents' characteristics. 
\newpage

\section{Game 1: Bilbo and the dragon}
This first experiment has been created with the aim of employing Q-Learning algorithm, Deep Q-Learning and Deep Q-Learning with the integration of Genetic Algorithm on a single system whose objective is beating a non-intelligent rival.\\*
The environment the three agents will be trained within is defined as a class named "World" in the Python programming language, by which it is possible to create the elements they interacts with during the learning process. First of all, it is initialized a Numpy array with dimensione 15x15 as an empty grid, to which various elements are added, so as to compose the world the agent explore:
\begin{itemize}
  \item Obstacles, whose presence in the grid is denoted by the special character '\ding{122}', and which can be situated randomly - if no "entrance" location is specified-, or in a precise position over the epochs;
  \item A treasure ($\symking$), which can be located in a randomly chosen coordinate or predefined too. It can not be moved, but only caught by the agent;
  \item A dragon ($\skull$), which is the non-intelligent enemy of the agent. It can be spawn randomly, or, in case the position of treasure was preselected, situated close to it, in order to make the game more difficult. In this experiment, the dragon is not capable to learn any strategy to win the game over the epochs or react to agent's moves, but it merely moves randomly through the grid following the four directions 'up', 'down', 'right', 'left';
  \item \textit{Bilbo} ($\smiley$), the agent to be trained, which is always generated at a random position. Unlike the other components, Bilbo moves around the world trying to maximize the total reward of his game and avoid penalties in best way possible, which is the core operating principle of any Reinforcement Learning methodology.
\end{itemize}
The creation of these components is carried out by the function \textit{"random\_spawn"} at locations where there are no obstacles or other characters, so that the grid remains consistent, while the function \textit{"is\_border"} checks whether a cell at a given coordinate is borderline.\\*
Therefore, Bilbo's goal is to reach the cell where the treasure is without getting caught by the dragon, condition that occurs when the two end up in the same cell. The game ends when Bilbo is killed by the dragon in the terms just described, or when he manages to get the treasure, escaping from the enemy among the obstacles scattered around the world. In the meanwhile, another function returns the reward based on every game state:
\begin{itemize}
%Sarebbe bello spiegare perchè avete scelto proprio questi valori:
  \item If the treasure has been caught by Bilbo, he gets a positive reward of 10;
  \item If Bilbo has been killed by the dragon, he gets a penalty of -5;
  \item At each moves, Bilbo receives a penalty of -1. This negative reward serves to incentivize Bilbo to find shorter routes to reach the treasure;
  \item If Bilbo dumps into an obstacle, so his position remains the same, he gets a penalty of -1, so that he learns to avoid them.
\end{itemize}
Reward values can be a negative constant or 0, but in order to obtain a faster convergence, it's better to give the agent a positive reward when he gets near the objective.\\*
Another Python class named "Agent" serves to set the main characteristics of Bilbo:
\begin{itemize}
  \item He can move in the same four direction of the dragon - 'up', 'down', 'right', 'left';
  \item He is affected by a "fear" factor, in order to simulate an aspect of real human exploring process. This is produced by a function which returns a value based on the distance from the dragon: if Bilbo is far from it, the "fear" value is close to 1; otherwise, as the distance from the dragon decrease, the value increases at most to 2. This measure is calculated as follows: 
  $$\frac{d}{1+d}$$
  with d representing the Euclidean distance between the dragon's position (D) and Bilbo's location (B):
  $$d=\sqrt{(D_{x}-B_{x})^2+(D_{y}-B_{y})^2}$$
  This will consist in a divisor of the $\epsilon$ parameter used in Reinforcement Learning algorithms to produce random moves in the grid and allow a better exploration of the world. As one can see from the formulation, this factor is a normalized distance, used to prevent $\epsilon$ to be too much affected by it.\\*
  In conclusion, in this way, when Bilbo approaches the dragon, his fear level increases and the chance of a random exploration as well.
\end{itemize}
Given this general aspects of Game 1, three methodologies have been exploited to provide the agent with a learning process, each based on an agent generalization provided by a new Python class.

\subsection{Q-Learning algorithm}
The first approach consists in Q-Learning algorithm, which is based on the propagation of potential reward from the best possible actions in future states. Following this idea, each action in each state is related to a Q value, updated when new relevant information about the environment, derived from exploration, is made available. The Q-Learning algorithm utilizes the following updating rule:
$$Q(s,a)=Q(s,a)+\alpha(r+\gamma\;max\;Q(s',a') - Q(s,a))$$
with $Q(s,a)$ Q value of current state and $Q(s',a')$ Q value of next state; $r$ the reward obtained by taking action $a$ at the state $s$; $\gamma$ discounting factor regulating the impact of future rewards; $\alpha$ learning rate. Hence, the worth of an action in a certain state also depends on the discounted worth produced by the best action in the next state $s'$, and on the extent of the steps to the convergence we selected with parameter $alpha$. In other words, this policy makes agent estimate the goodness of an action in a certain state based on the cascaded, discounted reward from the next states. In this way, at state $s$, Bilbo won't choose the action providing the best rewards only in the current state, but the action which leads to the best total rewards over all the states, while exploring the world in order to discover the rewards or penalties of all action. As Q-values are updated, they are stored in a Q matrix, so that Bilbo can learn from past experiences. This exploration process, that consists merely in random moves around the "gridworld", is fostered by $\epsilon$ parameter, which depends on "fear factor" as aforementioned. More over, $\epsilon$ decays by being multiplied to a "decay factor" at each step, up to a minimum of 0.01.\\*
In conclusion, this process can be explained as follows: if a number extracted from a uniform between 0 and 1 is less than $\epsilon$ or the Q values of all possible actions in the current state are null, the system performs a random step. Otherwise, it performs the action related to the maximum Q value in that state.\\*
The values set for hyperparameters are the following:
%Sarebbe bello spiegare perchè avete scelto proprio questi valori, anche qui:
\begin{itemize}
  \item $\alpha$ = 0.5, so that Bilbo's "patience level" is perfectly balanced
  \item $\gamma$ = 0.8, making Bilbo a fairly far-sighted agent 
  \item $\epsilon$ = 0.2
  \item $decay\;factor$ = 0.99, yielding a slow decay
\end{itemize} 
Then, the limits of epochs and episodes executable are both set to 1000.\\*
During the iterations, the performances of Bilbo are tracked by counting the number of wins, cases where he manages to stay alive and get the treasure, and losses, when he ends up killed by the dragon.
%Qui ci vanno i risultati

\subsection{Deep Q-Learning approach}
In the previous methodology, Q-Learning algorithm, the experience acquired by Bilbo during the games was based on explicit tables, - Q matrix and reward matrix - holding the information discovered, so that he could memorize which actions to choose in any given state. This tables could end up to be far too large and unwieldy when the game is composed of a large number of states and potential actions, or the world becomes larger. Similarly, still maintaining a medium or small gridworld, this mechanism leads to a loss of efficiency in the learning process, slowing down the progress of the system - namely the algorithm - which constitutes the agent.\\*
The solution to this problem can be the application of a Deep Reinforcement Learning methodology, which consists in training a neural network to predict Q values for each action in a certain state, instead of having explicit tables.\\*
The input of this neural network, implemented within Q-learning algorithm, is not only the simple set of coordinate describing the location of the agent itself and of the dragon, which is the current state, but a reshaped form of it, created through the function \textit{deep\_normalized\_state}. This one provides the input "image" for the Deep Q-Learning algorithm by:
\begin{enumerate}
  \item computing the distance between Bilbo and the treasure and between Bilbo and the dragon by subtracting respective coordinates;
  \item checking whether the position of Bilbo is at the boarder  for all four sides - ordered as "up", "down", "right", "left" - and assigning 0 if the given direction is accessible or 1 if it is borderline.
\end{enumerate}
After this transformation, the current state that will form the input of the neural network will appear as an array containing the two distances and the aforementioned binary values. For example, if Bilbo were on the leftmost edge of the gridworld, the input neurons would receive the following elements:
$$
\{B_{(x,y)}-T_{(x,y)},\;B_{(x,y)}-D_{(x,y)},\;0,\;0,\;0,\;1\}
$$
The output are the Q values for each action in the given state, which will approach the results produced by the learning update previously utilized. Therefore, using the mean-squared error metrics, the loss or cost function for the neural network will be:
$$(r+\gamma\;max_{a'}\;Q'(s',a')-Q(s,a))^2$$
Starting from this, the neural network is built as a sequential model composed by:
\begin{itemize}
  \item the input layer, activated by a Rectified Linear Unit (ReLU) function, receiving the reshaped current state, namely the information about the current distance from the other entities and with respect to the border; 
  \item an hidden layer composed by 8 neurons activated by a ReLU function too;
  \item the linear activated output layer with one neuron for each possible action in the given state, which hence performs the linear summation of the inputs and the weights, with no additional function applied, providing the estimated Q values.
\end{itemize}
The model is compiled using the aforementioned mean-squared error loss function and the Adaptive Moment Estimation, an extension of Stochastic Gradient Descent method, as optimizer.\\*
As Bilbo explores the world, his experience is described into 5 values: current and next state, game state, reward and action performed. Since, among all the moves performed, the probability that one of them leads to Bilbo's victory - attainment the treasure - or to his defeat - capture by the dragon - is very low, two separate memories are established: one containing the unlikely events just defined, the other storing the rewards obtained in the normal phases of the game. Therefore, when the system has developed sufficient experience in the world, in order to calculate the exact Q-values for each action in the given state, at each epoch the model is trained on a \textit{mini-batch} consisting of the union of 16 elements randomly extracted respectively from both separate memories. While the system performs these steps, it simultaneously updates the Q-values. Finally, this process is exploited within the previously used Q-Learning algorithm, extending the classic updating rule and resulting in a highly efficient Deep Reinforcement Learning methodology.\\*
As before, performances, computed by the number of win and losses, are tracked over the 20000 episodes, which involve at most 500 epochs.\\*
The hyperparameters are set as follows:
\begin{itemize}
  \item $\gamma$ = 0.8, maintaining the same Bilbo's "foresight" as before
  \item $\epsilon$ = 0.5, increasing the fixed component of random exploration probability with respect to the previous section
  \item $decay\;factor$ = 0.9999, further slowing the decay
\end{itemize} 
%Risultati

\subsection{Genetic Algorithm applied to Deep Q-Learning algorithm}
The third methodology arises from the need to optimize the weights of the neurons during the training of the neural network.\\*
A solution is represented by the Genetic Algorithm, which simulates the evolution of an initial population towards the best possible solution, the typical mechanism of the natural selection. During this process, each candidate solution present in the initial population of randomly generated individuals has a set of properties, chromosomes or genotype, which can be mutated and altered. The algorithm performs an iterative method which evaluates the fitness of every individual within each generation - actually the value of the objective function in the optimization problem being solved -, stochastically selects  the more fit individuals and forms a the new generation by combining and eventually mutating their genotypes. The algorithm terminates when a defined number of generations keeps on producing the same solution, or a satisfactory fitness level has been reached for the population.\\*

%setbilbo 
%Risultati
\newpage

\section{Multi-agent Competition}
\subsection{Game 2: Duel between two agents}
%Risultati
\subsection{Game 3: Battle among five agents}
%Risultati

\section{Multi-agent Cooperation}
%Risultati

\section{Conclusions}
Deep Learning per il problema multi-agente ed altre possibili migliorie.
%Conclusions should summarize the central points made in the Discussion section, reinforcing for the reader the value and implications of the work. If the results were not definitive, specific future work that may be needed can be (briefly) described. The conclusions should never contain ``surprises''. Therefore, any conclusions should be based on observations and data already discussed. It is considered extremely bad form to introduce new data in the conclusions.

\bibliographystyle{IEEEtran}
\bibliography{references.bib}

\end{document}
